\chapter{Counting animals in Camera-Trap Images}\label{cap:counting}

In this chapter, we present a proposed baseline for the problem of counting the
number of animals of each species in a sequence of camera trap images. Our
solution, based on a heuristic strategy using bounding boxes combined with a
species classifier, was ranked first in the iWildCam 2021 challenge.

\section{Introduction}

Camera traps have been widely used for monitoring wildlife, collecting large
amounts of images worldwide. In the last decade, the computer vision community
has investigated diverse approaches to improve automated information extraction
from this data, speeding up data preprocessing and reducing or eliminating the
need for manual data processing. This allows ecologists to focus on their
ecological research questions instead of spending time on manually reviewing
images. One of the widely studied data extraction tasks is species
identification, which provides information that enables scientists to model
species richness, species occurrence, and species distribution
\citep{beery2021iwildcam}. However, solely identifying the species is not
sufficient for other types of modeling, such as those focused on species
abundance or density. For these cases, counting the number of individuals
captured across image sequences is required.

Counting the number of individuals in a single image can be done using an
object detector such as MegaDetector. However, most camera trap projects
capture more than one picture per capture event in the hope of obtaining a good
shot to identify the photographed animal. In this context, simply counting the
number of bounding boxes in the pictures can overestimate or underestimate the
actual number of individuals in the sequence of images, as the same
individual may appear in multiple images, or different individuals may appear
in different images of the same sequence. Thus, it is necessary to use more
sophisticated approaches, such as tracking individuals or re-identifying them
across the images in the sequence. However, this is a challenging task due to
temporal discontinuities, as the images are usually taken at intervals of one
second or more, and most public camera trap datasets do not include count
labels, making supervised training unfeasible.

Some efforts have been made to address the counting of individuals in camera
trap images. For instance, \citet{norouzzadeh2018automatically} used a
classifier to count the number of animals in the Snapshot Serengeti dataset. In
a subsequent work, \citet{norouzzadeh2021deep} employed a pre-trained object
detection model to count animals at the image level based on the number of
bounding boxes. An effort to incentivize the development of methods for
counting the number of individuals across camera trap image sequences was the
launch of iWildCam 2021 Challenge \citep{beery2021iwildcam}, which provided a
dataset for this problem. Its purpose was to encourage the development of
approaches for the scenarios where count labels are not available, promoting
alternative methods for counting animals, with labels being collected only for
the test set.

Given this context, in this chapter, we present a robust baseline for counting
animals in sequences of camera trap images. Our approach combines a heuristic
based on bounding box counts with an ensemble of species classifiers.

\section{Materials and Methods}

In this section we present the iWildCam 2021, the dataset used to validate our
counting aproach, and detail our heuristic strategy that serve as a strong
baseline for the problem of counting animals in camera trap images.

\subsection{{iWildcam} 2021 dataset}

The iWildCam 2021 dataset comprises three components: camera trap
images from the WCS dataset, citizen science images from iNaturalist, and
multispectral imagery for the camera trap locations. The camera trap data,
provided by the Wildlife Conservation Society (WCS) \citep{wcs_lila}, contains
263,528 images of 206 species from 414 locations across 12 countries worldwide.
The official training/test split is based on camera location, with the training
set containing 203,314 images from 323 locations and the test set containing
60,214 images from 91 locations. As most publicly available camera trap data
does not include count labels, the count labels for image bursts were collected
exclusively for the test set to encourage the development of methods that can
learn to count without explicit labels \citep{beery2021iwildcam}.

The dataset also includes a subset of the iNaturalist 2017-2019 competition
dataset \citep{horn2018inat}, with 13,051 additional images from 75 species.
Although these images typically present a different data distribution and are of
higher quality, they can still be valuable for improving classification.

Finally, the dataset includes raw remote sensing data for each camera location,
collected by the Landsat 8 satellite, allowing investigation into whether the
model performance can be improved using this kind of multimodal data. GPS
coordinates for each camera location are provided, with most of them obfuscated
for privacy and security reasons. Each location has been randomly adjusted to
lie within 1 km of the original location center.

\subsection{A heuristic baseline for counting animals in camera-trap
images}

Our proposed heuristic baseline consists of two components: an image burst
classifier based on an ensemble of EfficientNet-B2 \citep{tan2019efficientnet}
models and a counting heuristic based on the bounding boxes generated by
MegaDetectorV4 \citep{beery2019efficient}. The main idea is to estimate the
number of animals present in the sequence based on the image with the highest
number of bounding boxes, which could provide a lower bound on the number of
animals across the sequence if the animal detector achieves perfect detection.
Additionally, since the vast majority of the images contain only a single
species, the simplest approach for classification is to assume a single species
across the entire sequence and use the most likely species prediction to the
total count of individuals.

Initially, MegaDetectorV4 is run on all images to detect animals. Next, we run
the image classifiers on both the full image and the bounding box with the
highest confidence score. The image prediction is calculated as the weighted
average of the predictions from both models. To predict the species in the
sequence, we average the species predictions of the non-empty images in the
burst. Finally, we generate a prediction vector containing the number of
individuals of each species -- in our case, only a single species -- based on
the maximum number of bounding boxes with confidence score higher than $0.8$
across any image in the sequence. If the classifiers identify only empty images,
the vector will contain only zeros. \autoref{fig:wildcam2021_winning}
presents a diagram summarizing the steps of our proposed heuristic baseline.

This classification strategy limits us to predicting only one species per
sequence. However, it has proven to be a strong baseline, as it was the winning
solution in iWildCam 2021. In the next subsections, we detail our
implementation.

\subsection{Species classifier}

We trained two species classifiers based on the EfficientNet-B2 architecture:
one trained using full images and another (bbox) trained with square crops containing
animals, extracted from bounding boxes generated by MegaDetectorV4. We chose
this approach following the iWildCam 2020 solutions, where the bbox model can
provide better predictions due to the highlighted animal. On the other hand, the
full image model has been shown to be more effective at classifying images with
animal herds. We used only the provided camera-trap data to train the classifier
models.

To train the bbox model, we treated all bounding boxes with a
confidence score equal to or higher than 0.6 as independent training samples and
applied square crops around them, using the size of the largerst bbox side. If
no bounding box met this confidence threshold, we used the full image for that
training instance. The image label was assigned as the bounding box label.

During inference, the species prediction assigned to each image is the weighted
average of the predictions from the bounding box with the highest score, the
full image, and their horizontally mirrored versions, following the approach
used by the winning solution of iWildCam 2020: 0.15 (full image) + 0.15
(mirrored full image) + 0.35 (bbox) + 0.35 (mirrored bbox). To classify a
sequence, we average the predictions of all nonempty images within the sequence.

\textbf{Handling class imbalance}. To address the class imbalance problem
inherent in real-world applications, such as the one represented in iWildCam
2021 dataset, we applied the Balanced Group Softmax (BAGS) method. Following the
original BAGS papers, we grouped classes into 4 softmax groups based to the
number $N$ of training instances: $N < 10$, $10 \le N < 100$, $100 \le N <
1000$, $N \ge 1000$. For each softmax group, we included the class ``others'' to
represent instances from all classes not included in that group. During
training, for each softmax, we undersampled the "others" category to ensure it
contained at most eight times the number of instances belonging to other classes
per batch, to prevent it from dominating the softmax. Additionally, we included
a special softmax group for the foreground/background classification. For the
final prediction, we remapped all predictions to the original softmax, ignoring
the ``others'' class. As in the original paper, these predictions are not true
probabilities since they do not sum up to one, but we consider the highest value
as the BAGS prediction.

\textbf{Image preprocessing}. For image preprocessing, we applied a square crop
centered on the bounding boxes or a random crop of aspect ratio
sampled in $[3/4, 4/3]$ and area in $[65\%, 100\%]$ for full image, followed
by a random horizontal flip and RandAugment (N=6, M=2). Then, the images were
normalized and resized to match the input dimensions of the network. During the
fix train/test resolution stage, we used inference preprocessing, which consists
only of normalizing and resizing the image or square crop to the input
dimensions of the network.

\textbf{Multi-stage training}. We trained the models in multiple stages. First,
each model was trained using the standard softmax with the default input
resolution of EfficientNet-B2 (260 x 260). In the first stage, only the
classifier layer was trained for 4 epochs, while the backbone -- previously
initialized with ImageNet weights -- remained frozen. In the second stage, all
layers were unfrozen, and all weights were fine-tuned for 20 epochs. In the
third stage, to fix the train/test resolution discrepancy issue
\citep{touvron2019fixing}, we fine-tuned only the last inverted bottleneck block
of the EfficientNet-B2 architecture and the classifier layer for 2 more epochs.
This was done using a higher input resolution (380 x 380) and applying the
inference preprocessing steps to the images. Afterward, we removed the plain
softmax and trained the BAGS header on top of the backbone, which was kept
frozen. In the fourth stage, the BAGS header was trained for 12 epochs using
training preprocessing with data augmentation at the inference resolution (380 x
380). Finally, in the fifth stage, the BAGS header was fine-tuned for 2 epochs
using inference-time preprocessing.
\autoref{table:countig_baseline_train_stages} presents a summary of the
parameters used.

\begin{table}[htb]
\caption{Training hyperparameters for each stage of the classifiers training.}
\label{table:countig_baseline_train_stages}
\begin{center}
\begin{tabular}{llllll}
\hline
 & \textbf{Weights trained} & \textbf{Epochs} &
\textbf{LR} & \textbf{Preprocessing} & \textbf{Resolution} \\
\hline
Stage 1* & Classifier layer                   & 4  & 0.01  & Training  & 260 x
260 \\
Stage 2  & All                                & 20 & 0.01  & Training  & 260 x
260 \\
Stage 3  & Last block and classif. & 2  & 0.001 & Inference & 380 x
380 \\
Stage 4  & BAGS header             & 12 & 0.01  & Training  & 380 x
380 \\
Stage 5  & BAGS heade             & 2  & 0.001 & Inference & 380 x
380 \\ \hline
\end{tabular}
\begin{tablenotes}
\footnotesize
\item [*] * For the bounding box model only.
\end{tablenotes}
\end{center}
\end{table}

\textbf{Implementation details}. To adjust the training hyperparameters, we
held out a validation set from the training set based on the locations.
However, to train the final models, we used all images available in the
original training set. The models were initialized with ImageNet weights and
trained with a batch size of $32$ using SGD with an initial learning rate of
$0.01$ ($0.001$ during the fix train/test resolution stages) and momentum of
$0.9$. The learning rate was linearly warmed up from $0$ to the initial value
over one-third of the steps in the first epoch of each training stage, and then
decayed to $0$ following the cosine schedule. We also applied label smoothing
with a value of $0.1$.

\section{Results}

The iWildCam 2021 challenge used the mean columnwise root-mean-squared error
(MCRMSE), described in \autoref{sec:count_metrics}, to evaluate solutions for
counting animals in camera traps. \autoref{table:iwildcam2021_results} presents
the final standings based on the private score, which was calculated using 50\%
of the test set. Our approach was ranked first in the challenge out of 42 teams
that entered the competition.

\begin{table}[htb]
\caption{The iWildCam 2021 final results on the private test set.}
\label{table:iwildcam2021_results}
\begin{center}
\begin{tabular}{ll}
\hline
\textbf{Solution}    & \textbf{Private score} \\ \hline
\textbf{1st place (ours)}        & \textbf{0.0293278}     \\
2nd place             & 0.0293786              \\
3rd place             & 0.0308570              \\
4th place             & 0.0308756              \\
5th place             & 0.0336513              \\
6th place             & 0.0339028              \\
Baseline - iWildCam 2020 Winner  & 0.0348171              \\
7th place             & 0.0376347              \\
8th place             & 0.0383956              \\
Baseline - All zeros             & 0.0385147              \\ \hline
\end{tabular}
\end{center}
\end{table}

As stated by the competition organizers \citep{beery2021iwildcam}, the MCRMSE
metric tends to be a small number even when the error in counts is large. This
occurs because camera traps typically have a small number of individuals and
because the model is double penalized for both species and counting errors. The
organizers provided some simple baselines, which we specify in
\autoref{table:iwildcam2021_results} for comparison purposes. One baseline
uses the iWildCam 2020 winner's species predictions combined with the maximum
number of bounding boxes with confidence score higher than 0.8 across the
sequence -- in line with our solution. The all zeros baseline predicts zero for
all instances, which performs surprisely well. This highlights how challenging
this counting problem is, given its inherent conditions.

We present an ablation study to demonstrate the effectiveness of certain design
choices in our heuristic solution, as shown in
\autoref{table:iwildcam2021_baseline_ablation}. This study considers only the
species classifier at the image level used for the final submission. A
well-known approach to improving classification predictions is using test-time
augmentation; in our solution, we added a horizontal flip to both full image and
bbox models. We also found that averaging the species predictions across images
in a sequence provides better performance than the majority voting approach.
Finally, we used a threshold of $0.8$ for counting MegaDetectorV4 bounding
boxes, as proposed by the organizers. We initially used $0.9$ during the
competition but found that it discarded too many valid bounding boxes.

\begin{table}[htb]
\caption{An ablation study of our baseline solution. The bold line represents
our final solution.}
\label{table:iwildcam2021_baseline_ablation}
\begin{center}
\begin{tabular}{llll}
\hline
\textbf{Test-time aug.} & \textbf{Seq. aggreg.} &
\textbf{MDv4 threshold} & \textbf{Private score} \\ \hline
-               & Voting  & 0.9 & 0.0314948 \\
Horizontal flip & Voting  & 0.9 & 0.0305960 \\
Horizontal flip & Average & 0.9 & 0.0294853 \\
\textbf{Horizontal flip} & \textbf{Average}  & \textbf{0.8} & \textbf{0.0293278}
 \\
\hline
\end{tabular}
\end{center}
\end{table}

\section{Final remarks}

Counting animals in sequences of camera trap images is a challenging problem,
especially considering the scenario where there is no count labels, as is
the case of most of public datasets. In this chapter, we presented a solution
based on a heuristic strategy that uses the bounding boxes generated by
MegaDetectorV4 combined with an ensemble of classifiers. Although our
strategy is limited to predicting only one species per sequence, it has been
proven to be a strong baseline for the problem, winning the iWildCam 2021
competition.

However, we believe that a more natural mechanism should be based on tracking
animals across images (multi-object tracking), classifying each track, and
counting them. Such an approach could also be more interpretable to humans,
avoiding ``black box'' solutions like those proposed by
\citet{norouzzadeh2018automatically}. During the competition, we experimented
with DeepSORT \citep{Wojke2017simple} and variations using embedding features
from ReID models. However, due to time constraints, we were unable to explore
this strategy in depth.

Our next steps on the counting animals problem include improving the proposed
baseline solution using the strategies explored in the other chapters, such as a
better empty-frame detection, long-tail approaches adapted to camera-trap data,
and more powerful classification models.

In the next chapter, we highlight the main contributions this work has already made to research in automated camera-trap image processing and outline future directions and next steps.